\subsection{Private Dissimilarity Estimation}

We are now ready to describe each mechanism in MTSP-LDP.
Recall that MTSP-LDP follows the dissimilarity guided publication framework,
where the server checks at every timestamp whether it is more beneficial to
approximate the current statistics with the last released statistics, than to
publish new statistics with noise.
This requires the server to compute the dissimilarity between user data $D_t$ at
time $t$ and the last released statistics.
It will be clear later that the last released statistics is also represented as
a private binary tree, denoted by $\OT_{t-1}$ (see \cref{ss:tree_publication}).
To compare the dissimilarity between $D_t$ and $\OT_{t-1}$, we can convert $D_t$
to a private binary tree $\DT_t$ using the method described in \cref{ss:tree},
and the dissimilarity is defined as the mean of the squared differences between
corresponding node properties in these two trees, i.e.,
\[
  \dis_t = \mathrm{dissimilarity}(\DT_t, \OT_{t-1}) \triangleq
  \frac{1}{|\DT_t|}\sum_{a\in \DT_t}(\DT_t[a] - \OT_{t-1}[a])^2,
\]
where $|\DT_t|$ denotes the number of nodes in the tree.

However, in the LDP setting, the server cannot directly access user data $D_t$,
hence $\DT_t$ is unknown.
Instead, the server can estimate $\DT_t$ by $\hDT_t$ using several FOs, as
explained in \cref{ss:tree}.
Next, we present an unbiased estimator of $\dis_t$.

\begin{theorem}\label{thm:dis}
  If $\hDT_t$ is estimated with privacy budget $\epsilon_{t,1}$, then $\hdis_t$
  is an unbiased estimate of $\dis_t$ and satisfies $\epsilon_{t,1}$-LDP, where
  \[
    \hdis_t\triangleq\frac{1}{|\hDT_t|}
    \sum_{a\in\hDT_t} (\hDT_t[a] - \OT_{t-1}[a])^2 - \var(\epsilon_{t,1}).
  \]
\end{theorem}

% \begin{proof}
%   The proof can be found in \cref{appendix:dis}.
% \end{proof}

\begin{proof}
We begin by computing the expectation of $\hdis_t$:
\[
\E(\hdis_t)
  = \E\big[ \frac{1}{|\hDT_t|}
    \sum_{a\in\hDT_t}(\hDT_t[a]-\OT_{t-1}[a])^2 - \var(\epsilon_{t,1}) \big]
\]
Since $\var(\epsilon_{t,1})$ is a constant, it can be factored out of the expectation:
\begin{align*}
  \E(\hdis_t)
  &= \frac{1}{|\hDT_t|}\sum_{a\in \hDT_t} \E\big[
    (\hDT_t[a] - \OT_{t-1}[a])^2\big]  - \E\big[ \var(\epsilon_{t,1}) \big] \\
  &= \frac{1}{|\hDT_t|}\sum_{a\in \hDT_t} \big[
    (\DT_t[a] - \OT_{t-1}[a])^2 + \var(\hDT_t[a])\big] - \var(\epsilon_{t,1}) \\
  &= \frac{1}{|\hDT_t|}\sum_{a\in \hDT_t} \big[
    (\DT_t[a] - \OT_{t-1}[a])^2 + \var(\hDT_t[a]) - \var(\epsilon_{t,1}) \big]
\end{align*}

Since $\hDT_t[a]$ is generated by the FO mechanism with
budget $\epsilon_{t,1}$, we have $\var(\hDT_t[a])=\var(\epsilon_{t,1})$. Thus, the expectation becomes:
\[
  \E(\hdis_t)
  = \frac{1}{|\hDT_t|}\sum_{a\in \hDT_t} \big[
  (\DT_t[a] - \OT_{t-1}[a])^2 \big]
  = dis_t
\]

Finally, because differential privacy is immune to post-processing~\cite{Dwork2013},
estimator $\hdis_t$ is still $\epsilon_{t,1}$-LDP as long as $\hDT_t$ is
$\epsilon_{t,1}$-LDP.
\end{proof}

\paragraph{Remarks}
% \remarks.
The dissimilarity is estimated at every timestamp using a fixed privacy budget
$\epsilon_{t,1}=\epsilon/(2w)$, i.e., half of the total privacy budget
uniformly spent on each timestamp in the window.
Since differential privacy is immune to post-processing~\cite{Dwork2013},
estimator $\hdis_t$ is still $\epsilon_{t,1}$-LDP as long as $\hDT_t$ is
$\epsilon_{t,1}$-LDP.
